\scanpage{41}
\mspage{20}
2. La question de pleine fidélité.

Soient $K,K'$ deux extensions de type fini de $\mathbb Q$ --- est-il vrai que tout $\Pi_{\mathbb Q}$-homomorphisme $\Pi_{K'}-\Pi_K$ provient d'un homomorphisme de corps $K'\longrightarrow K$? On est ramené aussitôt au cas où --- une clôture algébrique $\overline{\mathbb Q}$ de $\mathbb Q$ étant choisie, d'où un $\Gamma_{\overline{\mathbb Q}/\mathbb Q}$ --- $K$ et $K'$ ont des sous-corps $k,k'$ (clôture alg. de $\mathbb Q$ dans $K$ resp. $K'$) isomorphes, avec des \deleted{iso} \inserted{plongt} $k,\deleted{k'}\longrightarrow\mathbb Q$ de $\unsure{\overline m}$ image, \unsure{d'où} que $E_K$ et $E_{K'}$ peuvent être considérés comme des extensions d'un même groupe $\Gamma=\Gamma_{\overline{\mathbb Q}/k}$ par $\pi_K$ resp. $\pi_{K'}$. La question est alors si \uline{tout} hom. de $\pi_{K'}$ dans $\pi_K$ qui commute à l'action de $\Gamma$, est induit par un hom. $K\hookrightarrow K'$. Pour construire ce dernier, il faudrait donc avoir une idée comment reconstruire $K,K'$ à partir des extensions $E_K,E_{K'}$, ou encore à partir
\margin{40}
